\documentclass[11pt,a4paper]{article}
\usepackage[margin=22mm]{geometry}
\usepackage[T1]{fontenc}
\usepackage{lmodern}
\usepackage{microtype}
\usepackage[hidelinks]{hyperref}
\usepackage{parskip}
\setlength{\parindent}{0pt}
\pagestyle{plain}

\newcommand{\sourceheading}[2]{%
  \section*{#1}
  \textbf{#2}\par\medskip
}

\begin{document}

\begin{center}
  {\Large\bfseries LN905 Reading into Writing Practice}\par
  \vspace{4pt}
  {\large Gender -- 3 August 2026}
\end{center}

\section*{Question}

\textbf{``Gender bias at work is sustained mainly by organisations rather than by individuals.'' To what extent do you agree?}

Use all three extracts to develop your response. The extracts below are exam-style adaptations of the cited open-access papers. Their wording has been condensed and paraphrased while preserving the authors' central claims, evidence and stated limitations.

\sourceheading{Extract 1}{Adapted from Stamarski and Son Hing (2015), ``Gender inequalities in the workplace: the effects of organizational structures, processes, practices, and decision makers' sexism''}

Gender inequality in organisations cannot be explained only by the prejudiced choices of a few individuals. Human-resource systems shape who is hired, trained, evaluated, promoted and paid, so routine organisational practices can distribute opportunities unevenly even when a policy appears neutral. At the same time, managers bring their own assumptions into these systems. Hostile sexism can produce openly negative judgements of women, while benevolent sexism can appear positive but still restrict women by treating them as naturally caring, dependent or unsuited to demanding roles.

The authors therefore propose a reciprocal model. Organisational structures, leadership, strategy and culture influence HR policies and the conditions in which decisions are made. These conditions can make biased judgements more likely, particularly when evaluation criteria are vague. Decisions made under those conditions then reproduce the unequal organisation. The process can also work through socialisation: employees observe who advances, what conduct is rewarded and whether discrimination is tolerated. Over time, those messages may alter what decision makers regard as normal. A manager need not begin with strongly sexist attitudes to participate in a discriminatory system.

This model implies that individual awareness training is unlikely to be sufficient on its own. Organisations can reduce ambiguity by defining job-relevant criteria before evaluating candidates, using consistent measures of performance, providing decision makers with adequate information and time, and requiring them to justify their decisions. Changes to culture, leadership and accountability are also needed because isolated reforms may be weakened by the rest of the system.

The review does not establish that organisational structures always cause an observed disadvantage. Its evidence comes largely from Western, individualistic societies and from sectors dominated by men, and it focuses mainly on discrimination against women. The model is therefore best treated as an explanation of interacting mechanisms rather than a universal estimate of how much inequality is caused by institutions or individuals.

\textit{Source: Frontiers in Psychology, 6:1400. DOI: \href{https://doi.org/10.3389/fpsyg.2015.01400}{10.3389/fpsyg.2015.01400}. CC BY.}

\newpage

\sourceheading{Extract 2}{Adapted from Zedlacher and Yanagida (2023), ``Gender biases in attributions of blame for workplace mistreatment''}

Gender stereotypes may influence apparently individual judgements, but the direction of the bias can depend on the situation. Zedlacher and Yanagida tested this by showing standardised workplace videos to 880 members of the Austrian workforce. Eight professional actors appeared in 32 clips. The researchers varied whether the perpetrator and target were women or men and whether the incident involved an angry insult or exclusion from a shared lunch. Participants then judged matters such as intent, blame and moral anger.

The results did not show one simple rule in which every woman was judged more negatively than every man. Women who insulted a colleague were blamed at a similar level to men. However, a woman who excluded a female colleague from lunch was judged more harshly than a man in the same role. Female targets were also evaluated differently depending on the perpetrator's gender, while combinations involving male targets did not produce the same pattern. The form of mistreatment was therefore as important as gender in determining whether bias appeared.

These findings suggest that stereotypes guide observers most strongly when conduct is ambiguous and violates expectations about gender roles. A woman behaving in a way that conflicts with expectations of warmth may attract a different judgement from a man performing the same act. This supports organisational training that addresses specific attribution errors rather than presenting bias as a single, uniform attitude. It also suggests that organisations should examine how third parties interpret complaints, because biased judgements may affect whether intervention occurs.

The evidence should nevertheless be treated cautiously. The study examined only two forms of mistreatment in a controlled setting. Some effects were small, the three-way interaction effects were very small, and responses to short videos may not generalise to longer workplace relationships. One actor's likeability also affected some ratings. The experiment identifies context-dependent judgement patterns; it does not show how much workplace inequality these patterns produce outside the study.

\textit{Source: Frontiers in Psychology, 14:1161735. DOI: \href{https://doi.org/10.3389/fpsyg.2023.1161735}{10.3389/fpsyg.2023.1161735}. CC BY.}

\newpage

\sourceheading{Extract 3}{Adapted from Gygax, Garnham and Doehren (2016), ``What do true gender ratios and stereotype norms really tell us?''}

Beliefs about which occupations are usually performed by women or men are often described as gender stereotypes. Gygax and colleagues examined an important complication: some beliefs about occupational gender ratios may be reasonably close to the ratios people observe in the world. Earlier research had collected estimates for hundreds of occupations and social roles. When these estimates were compared with administrative data, the overall relationship was strong, although respondents tended to avoid the most extreme estimates.

Accuracy about current representation does not make a belief harmless. If most engineers in a society are men, recognising that fact is different from claiming that engineering is naturally better suited to men. Nor does a correlation between perceived and actual ratios reveal which one caused the other. People may learn ratios through experience, but both their beliefs and the labour-market distribution may also arise from shared influences such as media, education and parenting.

The causal direction can therefore run in a circle. A widespread belief that mathematics is masculine may influence how children judge their own ability and which careers they consider. Those choices can alter the later gender composition of occupations. The resulting composition then appears to confirm the original belief. In this account, an individual's apparently accurate judgement reflects a social reality that earlier beliefs and institutions helped to produce.

The study consequently warns against two conclusions. It is unjustified to treat every estimate of occupational representation as an additional personal prejudice, but it is equally unjustified to treat present-day ratios as evidence of biological suitability or free choice. The data cannot determine how much responsibility belongs to individual cognition, social norms or organisational opportunity. Effective intervention may need to change both the cues people encounter and the structures that turn expectations into unequal outcomes.

\textit{Source: Frontiers in Psychology, 7:1036. DOI: \href{https://doi.org/10.3389/fpsyg.2016.01036}{10.3389/fpsyg.2016.01036}. CC BY.}

\end{document}

